% ============================================================================
% appendix.tex
% ----------------------------------------------------------------------------
% Drop-in replacement for the appendix in PokerRL Writeup/paper.tex.
% Copy everything from "\appendix" downward, replacing the existing A.1--A.5.
%
% Cross-references used by the main body:
%   \ref{app:rules}    -> A.1  Leduc Hold'em Rules
%   \ref{app:state}    -> A.2  State and Opponent-Statistics Representation
%   \ref{app:hyper}    -> A.3  Training Configuration
%   \ref{app:oppstats} -> A.4  Opponent Statistics
%   \ref{app:pool}     -> A.5  Agent Pool and Evaluation Protocol
% ============================================================================

\appendix

% ----------------------------------------------------------------------------
\section{Leduc Hold'em Rules}
\label{app:rules}

Leduc Hold'em is a simplified two-player poker variant. The deck consists of
six cards --- two Jacks, two Queens, and two Kings --- and at every decision
point a player chooses from the action set
$\mathcal{A} = \{\textsc{Call},\,\textsc{Raise},\,\textsc{Fold}\}$.

The game has two betting rounds: \emph{pre-flop} (before a community card
is dealt) and \emph{flop} (after the community card is revealed). A raise
costs 2 chips in the pre-flop round and 4 chips in the flop round; at most
two raises are allowed per round.

At the start of each hand, both players post a 1-chip ante (initial pot of
2 chips) and each receives one private card. Pre-flop betting begins; if
neither player folds, a single public board card is then dealt and flop
betting follows. The hand ends in one of two ways:
\begin{enumerate}
  \item \textbf{Fold:} the opposing player wins the entire pot.
  \item \textbf{Showdown:} both cards are revealed. A player whose private
        card matches the board card has a pair and wins. If neither player
        has a pair, the higher private card wins
        ($\textsc{K} > \textsc{Q} > \textsc{J}$).
\end{enumerate}
Each agent observes only the information legally available to that player:
their own private card, the (possibly hidden) board, the pot vector, the
current round and acting player, the number of raises in the current
round, and the public action history.

% ----------------------------------------------------------------------------
\section{State and Opponent-Statistics Representation}
\label{app:state}

\paragraph{Game-state encoding.}
Every state $s$ is encoded as a 15-dimensional vector
$\phi(s) \in \mathbb{R}^{15}$ from the perspective of the \emph{viewer}
(the seat whose value is being estimated). The same encoding is used by
both the base value network and the modulation head.

\begin{table}[h]
\centering
\caption{15-dimensional game-state encoding $\phi(s)$.}
\label{tab:state}
\begin{tabular}{lrl}
\toprule
\textbf{Component} & \textbf{Dims} & \textbf{Description} \\
\midrule
Private hand            & 3 & One-hot over $\{\textsc{J},\textsc{Q},\textsc{K}\}$. \\
Board card              & 4 & One-hot over $\{\textsc{J},\textsc{Q},\textsc{K},\textsc{None}\}$. \\
Pot (viewer-relative)   & 2 & $(\text{viewer chips},\;\text{opp.\ chips}) / 13$. \\
My-turn flag            & 1 & $1$ if it is the viewer's turn, else $0$. \\
Position                & 1 & Seat index of the viewer ($0$ or $1$). \\
Current round           & 1 & $0$ pre-flop, $1$ flop. \\
Terminal flag           & 1 & $1$ if the hand has ended, else $0$. \\
Has-pair flag           & 1 & $1$ if the private card matches the board. \\
Raises this round       & 1 & Normalised by the per-round raise cap ($2$). \\
\bottomrule
\end{tabular}
\end{table}

The pot normaliser $13$ corresponds to the maximum total chip commitment
attainable at showdown under the betting rules of Section~\ref{app:rules}.

\paragraph{Modulation-head input.}
The value modulation head receives, in addition to $\phi(s)$, a
7-dimensional opponent-statistics vector $o$ summarising the current
opponent's observed behaviour over the ongoing session
(Section~\ref{app:oppstats}). The full modulation-head input is therefore
$[\phi(s)\,\Vert\,o] \in \mathbb{R}^{22}$.

% ----------------------------------------------------------------------------
\section{Training Configuration}
\label{app:hyper}

\subsection{Base Value Network}

The base value network $V_{\text{base}}(s)$ is a two-hidden-layer MLP
$15 \to 64 \to 64 \to 1$ with ReLU activations, trained with self-play
and a TD(0) objective. During training, action selection follows the
Boltzmann distribution induced by the one-step lookahead values with
temperature $\tau = 1.0$; at evaluation time, action selection is
greedy. Hyperparameters are listed in Table~\ref{tab:base-hyper}.

\begin{table}[h]
\centering
\caption{Training hyperparameters for the base value network.}
\label{tab:base-hyper}
\begin{tabular}{lr}
\toprule
\textbf{Hyperparameter} & \textbf{Value} \\
\midrule
Episodes              & $200{,}000$         \\
Learning rate         & $1\times 10^{-4}$   \\
Optimiser             & Adam ($\beta_1{=}0.9,\;\beta_2{=}0.999$) \\
Loss                  & MSE                 \\
Batch size            & $32$ episodes       \\
Architecture          & $15 \to 64 \to 64 \to 1$ (ReLU) \\
Temperature $\tau$    & $1.0$               \\
\bottomrule
\end{tabular}
\end{table}

\paragraph{TD(0) target.}
For a per-player post-action state chain $(s_1, s_2, \dots, s_T)$ with
terminal reward $r$, the bootstrap target is
\begin{equation*}
y(s_t) =
\begin{cases}
r & t = T, \\
\operatorname{sg}\!\bigl(V_{\text{base}}(s_{t+1})\bigr) & \text{otherwise},
\end{cases}
\qquad
\mathcal{L} = \mathrm{MSE}\!\bigl(V_{\text{base}}(s_t),\,y(s_t)\bigr),
\end{equation*}
where $\operatorname{sg}$ denotes the stop-gradient operator. Successive
states in a chain are separated by two game time steps, since the opponent
acts between consecutive learner actions.

\subsection{Value Modulation Head}

The modulation head implements
$V(s,o) = V_{\text{base}}(s) + \Delta(s,o)$, where $\Delta$ is a small MLP
mapping the concatenated 22-dimensional input through two hidden layers of
32 units to a scalar residual. The output-layer weights are initialised to
$\mathcal{U}(-10^{-2},\,10^{-2})$ and biases to zero, so that training
begins close to the frozen base. Hyperparameters are listed in
Table~\ref{tab:mod-hyper}.

The modulation head is trained against the weighted opponent pool of
Section~\ref{app:pool} (CFR and the heuristic agent oversampled $3\times$)
in sessions of $n = 100$ hands. The seat of the learner alternates between
episodes ($\text{seat} = \text{episode} \bmod 2$) so that both first- and
second-acting decisions are observed. Gradient updates apply to the
modulation parameters only, with $V_{\text{base}}$ frozen at the
self-play-trained checkpoint described above. We report mean and standard
deviation across $3$ random seeds.

\begin{table}[h]
\centering
\caption{Training hyperparameters for the value modulation head
(frozen-base variant; reported as ``Modulated Value Net.\ (Ours)''
in the main text).}
\label{tab:mod-hyper}
\begin{tabular}{lr}
\toprule
\textbf{Hyperparameter} & \textbf{Value} \\
\midrule
Episodes              & $300{,}000$             \\
Learning rate         & $1\times 10^{-4}$       \\
Optimiser             & Adam                    \\
Loss                  & MSE on residual TD(0) targets \\
Modulation-head arch.\ & $22 \to 32 \to 32 \to 1$ (ReLU) \\
Output init.          & $\mathcal{U}(-10^{-2},10^{-2})$, bias $0$ \\
Batch size            & $32$ episodes           \\
Session length $n$    & $100$ hands             \\
Prior strength $S$    & $20$                    \\
Seeds                 & $3$                     \\
\bottomrule
\end{tabular}
\end{table}

\paragraph{Residual TD(0) target.}
At each learner post-action state $s_t$ with stats snapshot $o_t$:
\begin{equation*}
y_\Delta(s_t,o_t) =
\begin{cases}
r - V_{\text{base}}(s_T) & t = T, \\[3pt]
\bigl(V_{\text{base}}(s_{t+1}) + \Delta(s_{t+1},o_{t+1})\bigr) - V_{\text{base}}(s_t)
& \text{otherwise},
\end{cases}
\end{equation*}
\begin{equation*}
\mathcal{L}_{\text{TD}(0)} = \mathrm{MSE}\!\bigl(\Delta(s_t,o_t),\,\operatorname{sg}(y_\Delta(s_t,o_t))\bigr).
\end{equation*}
The bootstrap target on non-terminal transitions is the \emph{total} next-state
value $V_{\text{base}} + \Delta$, so the residual head learns to be consistent
with itself; only the modulation parameters receive gradient updates, while
$V_{\text{base}}$ remains frozen throughout.

\paragraph{Optional state-conditioned gate.}
The state-gated variant studied in our ablations augments the composite
value as $V(s,o) = V_{\text{base}}(s) + g(s,o)\,\Delta(s,o)$, where
$g: \mathbb{R}^{22} \to [0,1]$ is a $22 \to 16 \to 16 \to 1$ MLP followed
by a sigmoid and initialised with near-zero output weights. Conditioning
the gate jointly on state and stats was motivated by the observation that
opponent style has variable impact across states.

\subsection{Baseline Reinforcement-Learning Agents}

The REINFORCE, Actor--Critic (A2C), and DQN baselines reported in
Section~4 are trained against the same weighted opponent pool used to
train our modulation head, with an opponent re-sampled at the start of
each episode (CFR $\times 3$, heuristic $\times 3$, six rule-based agents
$\times 1$; see Section~\ref{app:pool}). All three baselines share a
compute budget of $200{,}000$ episodes with learning rate
$1 \times 10^{-4}$, batch size $32$, and the Adam optimiser. Method-specific
settings are summarised in Table~\ref{tab:baseline-hyper}.

\begin{table}[h]
\centering
\caption{Baseline-specific hyperparameters.}
\label{tab:baseline-hyper}
\begin{tabular}{lll}
\toprule
\textbf{Setting} & \textbf{REINFORCE} & \textbf{Actor--Critic / DQN} \\
\midrule
Discount $\gamma$           & $1.0$                   & $0.99$ \\
Entropy bonus coefficient   & $0.01$                  & $0.01$ (A2C only) \\
Value-loss coefficient      & ---                     & $0.5$ (A2C only) \\
Replay buffer size          & ---                     & $10{,}000$ (DQN) \\
Target-network update       & ---                     & every $500$ episodes (DQN) \\
Exploration                 & Boltzmann (policy)      & $\varepsilon$-greedy, $1.0 \to 0.05$ over $200{,}000$ episodes (DQN) \\
\bottomrule
\end{tabular}
\end{table}

% ----------------------------------------------------------------------------
\section{Opponent Statistics}
\label{app:oppstats}

The opponent-statistics vector $o \in \mathbb{R}^{7}$ comprises six
round-stratified behavioural rates and a confidence scalar. Each rate is
maintained online by the learner during a 100-hand session and is reported
as a Beta--Bernoulli posterior mean smoothed toward a pool-average prior.

\paragraph{Posterior-mean smoothing.}
For each rate, let $k$ denote the count of qualifying events
(e.g.\ folds) and $n$ the count of qualifying contexts
(e.g.\ all pre-flop actions). With prior strength $S=20$ and pool-mean
prior $\bar p$,
\begin{equation*}
\hat p \;=\; \frac{k + \alpha}{n + \alpha + \beta},
\qquad
\alpha = \bar p \cdot S,\;
\beta  = (1 - \bar p) \cdot S.
\end{equation*}
The pool means $\{\bar p\}$ are calibrated once, before training, by
playing $500$ hands against each pool opponent with a uniform-random
learner (so that every action context is covered) and averaging the
observed counts across opponents.

\paragraph{Feature definitions.}
\begin{description}
\item[Pre-flop fold rate.] Fraction of pre-flop actions in which the
opponent folded. High values indicate a passive or tight pre-flop range.

\item[Pre-flop raise rate.] Fraction of pre-flop actions that were
raises. High values indicate aggression or wide pre-flop ranges.

\item[Flop raise rate.] Fraction of flop actions that were raises. A
spike relative to the pre-flop raise rate is a strong ``opponent paired
the board'' signal.

\item[Pre-flop fold-to-raise.] Among pre-flop actions taken while facing
a raise, the fraction that were folds. Distinguishes opponents who fold
to pressure pre-flop from those who call down.

\item[Flop fold-to-raise.] The analogous rate for the flop round.
Combined with the flop raise rate this characterises post-flop defence
frequency.

\item[Raise-after-raise rate.] Among actions taken while facing a raise
and with a re-raise legal, the fraction that were raises. A maniac-style
aggressor scores near $1.0$; passive opponents score near $0.0$.

\item[Confidence.] $\rho = n_h / (n_h + S)$, where $n_h$ is the number
of hands observed against the current opponent within the session and
$S = 20$. The scalar grows from $0$ at the start of a session to
$\approx 0.83$ after $100$ hands and allows the modulation head to
discount the other six features during the early, statistically-noisy
hands of a session.
\end{description}

% ----------------------------------------------------------------------------
\section{Agent Pool and Evaluation Protocol}
\label{app:pool}

\subsection{Training Pool}

The modulation head is trained against a fixed pool of eight opponents
that are sampled per session under the weighting in
Table~\ref{tab:train-pool}. The pool is deliberately unbalanced in skill:
CFR is a tabular Nash-equilibrium solver and represents the strongest
possible opponent in Leduc Hold'em, while the heuristic agent is a
strong, carefully tuned hand-crafted strategy that approximates expert
human play. The remaining six agents are simple style-based archetypes
(tight-passive, loose-aggressive, etc.) that exhibit exploitable
tendencies by construction and are correspondingly easier to beat. We
oversample CFR and the heuristic agent by $3\times$ so that, on a
per-session basis, the modulation head spends roughly half of its
training budget against the two highest-quality opponents in the pool.
This concentrates learning on the matchups that drive the robustness
metric --- a modulation head that performs well against weak rule-based
agents but collapses against CFR is of limited practical value.

\begin{table}[h]
\centering
\caption{Training-pool opponents and per-session sampling weights.
The expected per-session frequency for each agent is its weight divided
by the total weight ($14.0$).}
\label{tab:train-pool}
\begin{tabular}{lcl}
\toprule
\textbf{Opponent} & \textbf{Weight} & \textbf{Description} \\
\midrule
CFR              & $3.0$ & Tabular CFR$^+$; Nash-equilibrium baseline. \\
Heuristic        & $3.0$ & Hand-crafted expert-style strategy. \\
Tight-Passive    & $1.0$ & ``Nit''; raises only made pairs, folds Js to pressure. \\
Tight-Aggressive & $1.0$ & Raises premium hands; folds marginals under pressure. \\
Loose-Passive    & $1.0$ & Calling station; never folds pre-flop. \\
Loose-Aggressive & $1.0$ & Bluff-raises Js pre-flop (40\%) and on the flop (35\%). \\
Maniac           & $1.0$ & Raises whenever legal; folds only J-high under pressure with raises capped. \\
Random           & $1.0$ & Uniform-random over legal actions. \\
\bottomrule
\end{tabular}
\end{table}

\subsection{Evaluation Pool}

For the main results (Section~4 and Figure~2) we score each agent in a
\emph{fixed-gauntlet} round-robin against a 13-agent evaluation pool.
Eight opponents are in-distribution --- the same agents the modulation
head observed during training --- and five are out-of-distribution
neural-network agents that the modulation head never trained against
(Table~\ref{tab:eval-pool}).

\begin{table}[h]
\centering
\caption{Evaluation pool. The out-of-distribution agents were trained
independently for this study and never appeared in any modulation-head
training session.}
\label{tab:eval-pool}
\begin{tabular}{lll}
\toprule
\textbf{Opponent} & \textbf{Distribution} & \textbf{Notes} \\
\midrule
CFR              & In-distribution  & Tabular CFR$^+$ (Nash baseline). \\
Heuristic        & In-distribution  & Rule-based. \\
Tight-Passive    & In-distribution  & Rule-based archetype. \\
Tight-Aggressive & In-distribution  & Rule-based archetype. \\
Loose-Passive    & In-distribution  & Rule-based archetype. \\
Loose-Aggressive & In-distribution  & Rule-based archetype. \\
Maniac           & In-distribution  & Rule-based archetype. \\
Random           & In-distribution  & Uniform-random. \\
\midrule
Opp.\ Encoder    & Out-of-distribution & Frozen base $+$ learned opponent encoder. \\
Adaptive Value   & Out-of-distribution & Value net with 19-dim input (state $+$ 4 hand-coded stats). \\
REINFORCE        & Out-of-distribution & Policy gradient ($200{,}000$ pool-training episodes). \\
Actor--Critic    & Out-of-distribution & A2C ($200{,}000$ pool-training episodes). \\
DQN              & Out-of-distribution & Double Q-learning with replay and target network. \\
\bottomrule
\end{tabular}
\end{table}

\subsection{Match Protocol and Metrics}

Each (study agent, pool agent) matchup is played for $10{,}000$ hands with
alternating seats: the learner sits at seat $0$ for odd-indexed hands and
seat $1$ for even-indexed hands. Hands are grouped into $100$-hand
sessions, and the opponent-statistics tracker is reset at the session
boundary so that cold-start behaviour is repeatedly probed.

\paragraph{Per-matchup metrics.}
Each matchup yields the learner's mean chips per hand together with a
$95\%$ confidence interval and a \emph{cold} vs \emph{warm} decomposition
(hands $0$--$19$ vs hands $20$--$99$ within each session), used to isolate
the contribution of opponent-statistics accumulation.

\paragraph{Pool-level metrics.}
For each study agent we report, across the 13 pool opponents (and
separately across the in-distribution and OOD subsets):
\begin{align*}
\mu       &= \tfrac{1}{|P|}\sum_{p\in P}\bar x_p, &
\sigma    &= \operatorname{stddev}_{p\in P}(\bar x_p), \\
\mathrm{robustness} &= \mu - 1.5\,\sigma, &
\mathrm{worst\text{-}case} &= \min_{p\in P}\bar x_p,
\end{align*}
where $\bar x_p$ is the mean chips per hand against opponent $p$. The
robustness score penalises high variance across opponents, rewarding
agents that perform consistently rather than those that exploit some
opponents strongly while collapsing against others.

For comparability with the main text's tables, the per-100-hand session
score is $100 \times \bar x_p$ (e.g.\ $+0.419$ chips per hand
$=$ $+41.9$ chips per $100$-hand session).

\paragraph{Checkpoint selection.}
During training we evaluate every $10{,}000$ episodes against the
in-distribution pool ($500$ rounds per opponent) and retain both the
final checkpoint and the checkpoint that maximises the in-training
robustness score, which is then used for the final fixed-gauntlet
evaluation.
